\section{The tail sign dichotomy}\label{sec:dichotomy}

The vectors $a + b$ and $a - b$ partition the energy of the pair $(a,b)$
into two complementary directions. On the tail $H^c$, at least one of these
directions must carry significant energy---but it may be that both do, or that
one is negligible. This dichotomy determines which branch of the proof
applies.

\begin{lemma}[Tail sign dichotomy]\label{lem:tail-dichotomy}
\privateleanname{Main\.UniformAB\.tail\_sign\_dichotomy}
Let $a, b \in \R^n$ be unit vectors with $\ip{a}{b} = 0$ and set
$H = H(A,a,b)$. Exactly one of the following holds:
\begin{enumerate}[label=(\roman*)]
\item \textbf{Diffuse:}
$\norm{\tail_H(a - b)}_2^2 \ge 16\, s_0(A)$ and
$\norm{\tail_H(a + b)}_2^2 \ge 16\, s_0(A)$.

\item \textbf{Concentrated:}
there exists $\tau \in \{+1, -1\}$ with
$\norm{\tail_H(a - \tau b)}_2 < \theta(A)/32$.
\end{enumerate}
\end{lemma}

\begin{proof}
If $\norm{\tail_H(a - b)}_2^2 < 16\, s_0$, take $\tau = 1$; if instead
$\norm{\tail_H(a + b)}_2^2 < 16\, s_0$, take $\tau = -1$. Since
$16\, s_0 = A^4/2^{16} = (\theta/32)^2$, the squared bound yields the norm
bound in~(ii). If neither holds, both squared norms are at least $16\, s_0$
and we are in~(i).
\end{proof}

\begin{remark}
In the concentrated case, $a$ and $\tau b$ nearly agree on the tail, so
their difference is concentrated on the finite head set $H$. In the diffuse
case, both $a \pm b$ have energy spread across many small tail coefficients,
precisely the regime where anticoncentration is effective.
\end{remark}
